% 表1 展示第三方模型给 rationale 进行质量打分
\begin{table*}[t]
\centering
\normalsize
\setlength{\tabcolsep}{7pt}
\renewcommand{\arraystretch}{1.3}
% ---------- 上半部分：对决区 ----------
\begin{tabularx}{\textwidth}{@{}X|X@{}}
\toprule
\multicolumn{2}{@{}p{\dimexpr\textwidth-2\tabcolsep-\arrayrulewidth\relax}@{}}{%
\textbf{Task: Grounding Specificity} \quad $\vert$ \quad
\textbf{Evaluation: Rationale-Quality Assessment} \hfill
\textbf{Gold label $=3$}} \\
\midrule
\multicolumn{2}{@{}>{\raggedright\arraybackslash}p{\dimexpr\textwidth-2\tabcolsep-\arrayrulewidth\relax}@{}}{%
\textbf{Evaluation input:} \textit{``Limitations of modeling with GP:
We know that the required number of samples to accurately model the
objective function grows exponentially, leading to increased computational
complexity and potentially reduced accuracy. Do you plan to address this
limitation in future work?''}} \\
\midrule
\rowcolor{gray!12}
\textbf{SO\textsubscript{RF} rationale} \hfill \textbf{predicts 1} &
\textbf{RS\textsubscript{RF} rationale} \hfill \textbf{predicts 3} \\
The answer refers to a general limitation of modeling with GP
(Gaussian Processes) but does not specifically identify a section, table,
figure, or unique aspect of the paper. It is a general comment about a
known limitation in the field rather than a specific issue within the
paper. The comment does not specify what needs to be addressed in the
paper's context, making it unspecific in terms of the paper's content. &
The comment refers to ``Limitations of modeling with GP'' as a general
topic, but it does not explicitly mention a specific section, table,
figure, or unique aspect of the paper. The authors can make an educated
guess that it relates to the limitations or discussion of GP modeling,
but they cannot confidently pinpoint the exact part being addressed.
However, the comment is specific: it clearly identifies the issue
(exponential growth in required samples leading to computational
complexity and reduced accuracy) and asks whether the limitation will
be addressed in future work. This aligns with score 3: weakly grounded
but specific. \\
\end{tabularx}
\nopagebreak
\vspace{-2pt}
% ---------- 下半部分：四维评分表 ----------
\begin{tabularx}{\textwidth}{@{}l
    >{\centering\arraybackslash}X
    >{\centering\arraybackslash}X
    >{\centering\arraybackslash}X
    >{\centering\arraybackslash}X@{}}
\midrule
 & \textbf{Rubric correctness} & \textbf{Evidence grounding} &
   \textbf{Decision coverage} & \textbf{Unsupported-claim control} \\
\textbf{Judge on SO} & $1.000$ & $1.000$ & $1.500$ & $1.000$ \\
\rowcolor{gray!12}
\textbf{Judge on RS} & $3.000$ & $3.000$ & $3.000$ & $3.000$ \\
\midrule
\multicolumn{5}{@{}p{\dimexpr\textwidth-2\tabcolsep\relax}@{}}{%
\hfill \textbf{Overall preference: RS.}} \\
\bottomrule
\end{tabularx}
\caption{A project sample for four-dimensional rationale-quality assessment.
Each sample pairs a task with an embedded evaluation question; SO and RS
denote score-only and rationale-and-score training under rationale-first
inference.}
\label{tab:rationale-quality-example}
\end{table*}

% 表2 展示第三方模型给 rationale 进行一致性评判
\begin{table*}[t]
\centering
\normalsize
\setlength{\tabcolsep}{7pt}
\renewcommand{\arraystretch}{1.3}
% ---------- 上半部分：对决区 ----------
\begin{tabularx}{\textwidth}{@{}X|X@{}}
\toprule
\multicolumn{2}{@{}p{\dimexpr\textwidth-2\tabcolsep-\arrayrulewidth\relax}@{}}{%
\textbf{Task: Helpfulness} \quad $\vert$ \quad
\textbf{Evaluation: Rationale--Score Consistency Assessment} \hfill
\textbf{Gold label $=5$}} \\
\midrule
\multicolumn{2}{@{}>{\raggedright\arraybackslash}p{\dimexpr\textwidth-2\tabcolsep-\arrayrulewidth\relax}@{}}{%
\textbf{Evaluation input:} \textit{``Overall training time. The overall
training time and the bandwidth of exchanging hidden vectors should be
mentioned. If the network size of the exchanged information is too high
and bandwidth is limited in wireless communication, a single iteration
will take too much time or fail because of communication-protocol
constraints. For example, in an IoT setting, the maximum payload length
is only 65,535 bytes. Please discuss the size of the exchanged information
(hidden vectors or gradients) and the total training time in the revision.''}} \\
\midrule
\rowcolor{gray!12}
\textbf{SO\textsubscript{RF} rationale} \hfill \textbf{predicts 5} &
\textbf{RS\textsubscript{RF} rationale} \hfill \textbf{predicts 4} \\
The answer provides a clear and actionable suggestion for the authors to
address the training-time and bandwidth issues in their paper. It
specifically points out the need to discuss the size of the exchanged
information and the total training time, which are directly relevant to
wireless communication and IoT settings. The comment also justifies why
this is important by citing constraints such as the maximum payload length
in IoT. The feedback is comprehensive and detailed, making it a strong
example of a high-quality review comment. &
The comment identifies a specific weakness: the paper does not mention
the overall training time or bandwidth required for exchanging hidden
vectors. It clearly explains why this matters: if the exchanged information
is large and bandwidth is limited, an iteration could be slow or fail under
protocol constraints. It also explicitly asks the authors to report the
size of the exchanged information and total training time. The feedback is
clear, constructive, and directly useful, but it could be expanded with
more detailed guidance on how to measure or report these quantities.
Therefore, it aligns with a score of 4. \\
\end{tabularx}
\nopagebreak
\vspace{-2pt}
% ---------- 下半部分：一致性判决表 ----------
\begin{tabularx}{\textwidth}{@{}l
    >{\centering\arraybackslash}p{0.22\textwidth}
    >{\raggedright\arraybackslash}X@{}}
\midrule
 & \textbf{Verdict} & \textbf{Judge explanation} \\
\textbf{Judge on SO} & \texttt{partially\_supported} &
The rationale accurately identifies the comment's strengths, but calling
it fully comprehensive overstates its depth. The feedback focuses on one
issue and does not provide detailed guidance on measurement, reporting,
or mitigation, so it only partially supports a score of 5. \\
\rowcolor{gray!12}
\textbf{Judge on RS} & \textbf{\texttt{supported}} &
The rationale distinguishes the comment's clear, actionable strengths from
its limited scope and correctly explains why these properties support a
score of 4 rather than the highest score. \\
\midrule
\multicolumn{3}{@{}p{\dimexpr\textwidth-2\tabcolsep\relax}@{}}{%
\hfill \textbf{Overall preference: \underline{RS}.}} \\
\bottomrule
\end{tabularx}
\caption{A project sample for rationale--score consistency assessment.
Each rationale is judged for whether its content supports the score it
predicts; SO and RS denote score-only and rationale-and-score training
under rationale-first inference.}
\label{tab:rationale-consistency-example}
\end{table*}

% 表3 展示对错误样本的 rationale 进行句子级别的标注情况
\begin{table*}[t]
\centering
\normalsize
\setlength{\tabcolsep}{7pt}
\renewcommand{\arraystretch}{1.3}
% ---------- 上半部分：任务与被评 rationale ----------
\begin{tabularx}{\textwidth}{@{}X@{}}
\toprule
\multicolumn{1}{@{}p{\dimexpr\textwidth-2\tabcolsep\relax}@{}}{%
\textbf{Task: Actionability} \quad $\vert$ \quad
\textbf{Evaluation: Sentence-Level Error Annotation} \hfill
\textbf{Gold label $=3$}} \\
\midrule
\multicolumn{1}{@{}p{\dimexpr\textwidth-2\tabcolsep\relax}@{}}{%
\textbf{SO\textsubscript{DS} prediction: 3 \,(\textcolor{green!50!black}{correct})}
\hfill
\textbf{SO\textsubscript{RF} prediction: 5 \,(\textcolor{red!70!black}{severe error})}} \\
\midrule
\multicolumn{1}{@{}>{\raggedright\arraybackslash}p{\dimexpr\textwidth-2\tabcolsep\relax}@{}}{%
\textbf{Evaluation input:} \textit{``The authors claimed
convolutional-coding-based DUNs can transmit high-capacity informative
features among stages, however the initial convolutional coefficients
$\alpha$ and the convolutional dictionary $D$ are all learned by an E2E
way, which means there is no additional information brought into the
network during training. The constrain that projects the original signal
into high-capacity feature-level image representation seems can also be
achieved in the image domain by more hidden layer channels (more
parameters). It does not have insights. I suggest authors can give more
convincing motivation (e.g., theorem, experiments) to explain why such
operations are used.''}} \\
\midrule
\rowcolor{gray!12}
\multicolumn{1}{@{}p{\dimexpr\textwidth-2\tabcolsep\relax}@{}}{%
\textbf{RF rationale under annotation}} \\
The answer provides a specific suggestion for the authors to ``give more
convincing motivation (e.g., theorem, experiments) to explain why such
operations are used.'' \textcolor{red!70!black}{This is an explicit action
with concrete details on how to implement it.}
\textcolor{red!70!black}{The authors can directly identify what needs to be
done and how to apply the action.} \textcolor{red!70!black}{Therefore, this
meets the criteria for a score of 5.} \\
\end{tabularx}
\nopagebreak
\vspace{-2pt}
% ---------- 下半部分：句级错误标注表 ----------
\begin{tabularx}{\textwidth}{@{}p{0.30\textwidth}
    >{\raggedright\arraybackslash}p{0.20\textwidth}
    >{\raggedright\arraybackslash}X@{}}
\midrule
\rowcolor{gray!12}
\textbf{Rationale sentence} & \textbf{Error type} &
\textbf{Judge explanation} \\
{This is an explicit action with concrete details
on how to implement it.} &
Rubric misapplication &
``theorem, experiments'' names evidence categories but does not specify how
to execute them. \\ \addlinespace[3pt]
{The authors can directly identify what needs to be
done and how to apply the action.} &
Unsupported inference &
The comment does not provide an experimental design or theorem statement. \\ \addlinespace[3pt]
{Therefore, this meets the criteria for a score
of 5.} &
Score-mapping error &
The rationale supports an explicit but vague action, corresponding to score
3 rather than score 5. \\
\midrule
\multicolumn{3}{@{}p{\dimexpr\textwidth-2\tabcolsep\relax}@{}}{%
\hfill \textbf{Consensus: supports\_wrong\_score.}} \\
\bottomrule
\end{tabularx}
\caption{A project sample for sentence-level rationale error annotation.
Each rationale sentence is independently judged, and the per-sentence
verdicts are aggregated into a consensus label for the whole rationale.
Erroneous sentences are highlighted in red in the rationale and echoed in
the left column of the annotation table.}
\label{tab:rationale-sentence-example}
\end{table*}

% 表4 纠正错误句子后的 rationale，以及分数的变化
\begin{table*}[t]
\centering
\normalsize
\setlength{\tabcolsep}{7pt}
\renewcommand{\arraystretch}{1.3}
% ---------- 上半部分：对决区 ----------
\begin{tabularx}{\textwidth}{@{}X|X@{}}
\toprule
\multicolumn{2}{@{}p{\dimexpr\textwidth-2\tabcolsep-\arrayrulewidth\relax}@{}}{%
\textbf{Task: Actionability} \quad $\vert$ \quad
\textbf{Evaluation: Rationale Correction \& Rescoring} \hfill
\textbf{Gold label $=2$}} \\
\midrule
\multicolumn{2}{@{}p{\dimexpr\textwidth-2\tabcolsep-\arrayrulewidth\relax}@{}}{%
\textbf{RS\textsubscript{DS} prediction: 5 \,(\textcolor{red!70!black}{severe error})}
\hfill
\textbf{RS\textsubscript{RF} prediction: 2 \,(\textcolor{green!50!black}{corrected})}} \\
\midrule
\multicolumn{2}{@{}>{\raggedright\arraybackslash}p{\dimexpr\textwidth-2\tabcolsep-\arrayrulewidth\relax}@{}}{%
\textbf{Evaluation input:} \textit{``Somewhat limited scope of
experiments. E.g.\ do the findings hold up on ImageNet?''}} \\
\midrule
\rowcolor{gray!12}
\multicolumn{1}{@{}p{\dimexpr 0.5\textwidth-1.5\tabcolsep-\arrayrulewidth\relax}@{}}{%
\textbf{Original rationale}} &
\multicolumn{1}{@{}p{\dimexpr 0.5\textwidth-1.5\tabcolsep-\arrayrulewidth\relax}@{}}{%
\textbf{Corrected rationale}} \\
The comment points out a limitation in the experimental scope and asks
\textcolor{red!70!black}{a specific question}: ``do the findings hold up
on ImageNet?'' \textcolor{red!70!black}{This is an explicit request for
the authors to test their findings on ImageNet. The action is directly
stated, and the implementation is concrete: the authors know exactly what
to do (run experiments on ImageNet) and what to evaluate (whether the
findings hold up). Therefore, the comment is fully actionable with
explicit and concrete details.} &
The comment points out a limitation in the experimental scope and asks
\textcolor{blue!70!black}{a question}: ``do the findings hold up on
ImageNet?'' \textcolor{blue!70!black}{The comment is phrased as a question
rather than an explicit request, so the action of testing the findings on
ImageNet must be inferred by the authors. Moreover, even if treated as an
explicit statement, the implied action is vague: it does not specify how
to run experiments on ImageNet (e.g., which baselines, settings, or
metrics to compare). Thus the comment lacks concrete details on how to
implement the action.} \\
\midrule
\multicolumn{2}{@{}p{\dimexpr\textwidth-2\tabcolsep-\arrayrulewidth\relax}@{}}{%
\textbf{External correction:} editor \texttt{glm-5.3-flash}, reviewer
\texttt{MiniMax-M3}, status \texttt{approved}. \hfill
\textbf{Error types:} evidence misread, rubric misapplication,
score-mapping error.} \\
\end{tabularx}
\nopagebreak
\vspace{-2pt}
% ---------- 下半部分：定前缀重打分表 ----------
\begin{tabularx}{\textwidth}{@{}l
    >{\centering\arraybackslash}X@{}}
\midrule
 & \textbf{Rescored prediction} \\
\textbf{A: natural rationale and score} & $5$ \\
\textbf{B: fixed original rationale and score} & $5$ \\
\rowcolor{gray!12}
\textbf{C: fixed corrected rationale and score} & \textbf{$2$} \\
\midrule
\multicolumn{2}{@{}p{\dimexpr\textwidth-2\tabcolsep\relax}@{}}{%
\hfill \textbf{Result: B reproduces A; C recovers the gold label.}} \\
\bottomrule
\end{tabularx}
\caption{A project sample for rationale correction, independent review,
and fixed-prefix rescoring. Red marks the original statements that were
revised, and blue marks the revised statements in the corrected rationale.
Conditions A--C rescore the model with the natural, fixed original, and
fixed corrected rationale prefixes, respectively.}
\label{tab:rationale-correction-example}
\end{table*}
